\documentclass[11pt,a4paper]{article}
\usepackage[utf8]{inputenc}
\usepackage[T1]{fontenc}
\usepackage[a4paper,margin=24mm]{geometry}
\usepackage{amsmath,amssymb,mathtools}
\usepackage{booktabs,tabularx,enumitem}
\usepackage[expansion=false]{microtype}
\usepackage[colorlinks=true,linkcolor=blue,citecolor=blue,urlcolor=blue]{hyperref}
\setlist{itemsep=2pt,topsep=3pt}
\makeatletter
\g@addto@macro\UrlBreaks{\do\_\do\.\do\A\do\B\do\C\do\D\do\E\do\F\do\G\do\H\do\I\do\J\do\K\do\L\do\M\do\N\do\O\do\P\do\Q\do\R\do\S\do\T\do\U\do\V\do\W\do\X\do\Y\do\Z}
\setlength{\emergencystretch}{2em}
\makeatother
\newcommand{\lean}[1]{\nolinkurl{#1}}
\newcommand{\Rea}{\operatorname{Re}}
\newcommand{\NK}{\mathrm N}
\newcommand{\OO}{\mathcal O}

\title{What the dead code says:\\ abandoned routes in the Lean proof of the quasi-Riemann hypothesis}
\author{Notes on the pruning of \texttt{openai/math}, \texttt{lean/OAI/NumberTheory/DirichletL}}
\date{October 9, 2026}

\begin{document}
\maketitle

\begin{abstract}
The Lean development behind OpenAI's quasi-Riemann hypothesis (zero-free half-plane
$\Rea s>7/8$ for Dirichlet and Hecke $L$-functions, and the shorter $23/24$ proof) contains about
$6{,}400$ named declarations, roughly $99{,}000$ source lines, that no Comparator-checked result
uses. We pruned them. This note reconstructs, from the deleted code alone, what those abandoned
branches were trying to do and why they were dropped. It then lists the few ideas that look worth
more compute and the dead ends worth avoiding. Almost all of the deleted material is \emph{finished}
mathematics: there is no \lean{sorry} anywhere. What it records is a search process: interfaces
guessed before their inputs existed, theorems re-proved with a different quantifier order, and a
handful of genuine obstructions.
\end{abstract}

\section{Method and caveats}

\paragraph{Data.} Reachability was computed inside Lean, through types and proof terms, from the
results that must survive:
\begin{itemize}
\item the 7/8 theorems for $\zeta$, Dirichlet and Hecke $L$-functions;
\item the new 23/24 theorems for Dirichlet $L$-functions and $\zeta$;
\item the one module imported from outside the directory.
\end{itemize}
A declaration counts as \emph{dead} if it is unreachable; Table~\ref{tab:areas} gives the counts by
area. Five agents each audited one area. They read the original sources at commit \texttt{fd4aeeb},
diffed dead modules against their live successors, and searched for hypotheses that nothing ever
discharges. This note condenses their reports.

\paragraph{Caveats.} The repository has no development history: it was published as a single
commit. ``Why it was abandoned'' is therefore an inference from three kinds of evidence:
\begin{itemize}
\item \emph{the diff between a dead theorem and its live twin}, usually a change of one hypothesis
or one quantifier;
\item \emph{hypotheses that nothing constructs}: a structure or Prop that only ever appears as an
assumption;
\item \emph{import depth}, used as a proxy for chronology.
\end{itemize}
We mark conjectures as such. ``Dead'' means unused by the final theorems, not wrong: every dead
declaration compiles.

\begin{table}[h]
\centering\small
\begin{tabular}{@{}lrrr@{}}
\toprule
Area (directories) & dead declarations & lines & files removed \\
\midrule
Moments (7/8 fourth moments) & 1,480 & 24,235 & 95 \\
Descent, Inversion, Dictionary (7/8) & 952 & 17,724 & 76 \\
Theta function, reflection, cusps, Hecke & 1,230 & 16,008 & 49 \\
Sieves, Gauss sums, mean square & 1,200 & 20,301 & 0 \\
Detector, energy, Mellin step, extraction & 1,557 & 20,711 & 33 \\
\midrule
Total & 6,419 & 98,979 & 253 \\
\bottomrule
\end{tabular}
\caption{Unreachable named declarations by area, and the source lines they span. Files are removed
when nothing in them survives.}\label{tab:areas}
\end{table}

\section{Why code died: a taxonomy}

\begin{enumerate}[label=\textbf{(\Alph*)}]
\item \textbf{Superseded copies; the majority.} A module was cloned, one hypothesis or quantifier
changed, and the whole chain proved again. Some examples:
\begin{itemize}
\item In \lean{Descent/} and \lean{Inversion/}, 39 of 76 deleted modules have a surviving copy
that shares at least $70\%$ of its lines.
\item In \lean{Reflection/}, about $2.8$k of $5.2$k dead lines are twins that differ only in where
an $\exists$\,\lean{degree} sits.
\item In \lean{Moments/}, the chain for the $\Theta$-rows was rebuilt three times. The last rebuild
only removed the hypothesis \lean{IsCoprime Q C}, paying instead a constant factor
$\NK(\mathrm{rad}\,Q)^{2/3}$.
\end{itemize}
\item \textbf{Conditional chains whose hypothesis was never discharged.}
\begin{itemize}
\item \lean{FixedTestPrimeMS}: the prime-row extraction of the 11/12 paper. It needs many prime
ideals in a norm window.
\item \lean{CanonicalLowerMoments}: a strong induction over a real parameter, which is not
well-founded as stated.
\item A single-cusp reflection hypothesis \lean{hreflect}.
\item Per-row ramification certificates.
\item A Mertens-type prefix bound.
\end{itemize}
In every case the live proof won by \emph{eliminating} the hypothesis, not by proving it.
\item \textbf{Unneeded strength.}
\begin{itemize}
\item Contour estimates uniform in every parameter, where the paper says only the power margins
need to be uniform.
\item A full spectral theory of the Laplacian: self-adjointness, global meromorphy, a
finite-dimensional theta space.
\item Uniformity over a parametric family of test functions, where explicit seminorm products were
needed.
\end{itemize}
\item \textbf{Genuine mathematical dead ends.} These are few, and are listed in
Section~\ref{sec:deadends}.
\item \textbf{Live by replacement.} Upstream's own unconditional 23/24 assembly is dead only
because the new paper-numbered top level re-proves it from the same lemmas. It was never abandoned.
\end{enumerate}

\section{Findings by area}

\subsection{The cubic theta function and its reflection}
Paper~2 (the 11/12 paper) reduces the reflection at a cusp $a/c$ to three standard cusps. It does
so by asserting, without proof, that Kubota's $\theta$ is invariant under $\mathrm{SL}_2(\mathbb
Z)$. The dead code shows the formalizers meeting exactly this point.
\begin{itemize}
\item \textbf{The early chain could not get past the principal cusp.} It used only the Kubota group
$\Gamma(3)$ (\lean{levelThree}) and stopped there: the arithmetic structure requires $3\mid c_0$.
The Weyl-type cusps ($(c,\lambda)=1$) need the $\mathrm{SL}_2(\mathbb Z)$-invariance.
\item \textbf{The workaround was to average.} The live proof averages $\theta$ over the 24 cosets of
$\Gamma(3)$ in $\Gamma(3)\mathrm{SL}_2(\mathbb Z)$. The average is automorphic by construction, and
its Fourier coefficients are an explicit mixture $\tfrac14(\theta_h+3\theta^{\rm opp}_h)$ of the
expansions at $\infty$ and at the opposite cusp.
\item \textbf{The invariance itself is never stated or proved.} No dead or live declaration asserts
it, so the formal proof is correct without it.
\end{itemize}

Two further routes died:
\begin{itemize}
\item \textbf{The spectral and Green-identity routes to the residue.} The Green route stalled before
any Maass--Selberg relation on a truncated fundamental domain. The live residue comes from a
Fredholm pencil with an explicit Laurent split at $s=4/3$; only that local interface is used.
\item \textbf{An Euler-factor recurrence for $\tau(hp^3)=|p|\tau(h)$.} It solves to $A+Bq^k$, which
leaves one constant free, so it cannot prove the relation. The live proof shows instead that the
difference is a decaying $8/9$-eigenfunction orthogonal to the residue, hence zero.
\end{itemize}

\subsection{The top level of the 23/24 proof: Mellin step and extraction}
\begin{itemize}
\item \textbf{Six interfaces for the Mellin step.} The step ``power saving $\Rightarrow$ zero-free
half-plane'' went through six successive interfaces:
\begin{enumerate}
\item an exponentially smoothed rational M\"obius sum;
\item the same sum over ideals of $\mathbb Z[\omega]$;
\item sharp-cutoff prefix sums;
\item compactly supported smooth weights;
\item the same with the Euler factors at $S$ deleted;
\item prime rows.
\end{enumerate}
Each is a complete reduction to a hypothesis. Each was dropped because the one analytic input ever
proved, the mean-square bound, controls only smooth, compactly supported sums over ideals prime to
$S$. The interfaces were written about 120 import levels before that input existed.
\item \textbf{Paper~2's prime-row extraction was formalized but never closed.} It is conditional on
\lean{FixedTestPrimeMS}, which asks for $\gg Y^{1-\varepsilon}$ prime ideals of norm about
$Y=D^{11/60}$, and nothing constructs it. The integer rows $b^6$, with a totient weight, replaced
it. They give \emph{exactly the same exponent} $11/12+5\vartheta/12$ and need no prime ideal
theorem.
\item \textbf{Rational-prime rows are strictly worse.} An earlier variant used rows $r^6$ for
rational primes $r$. It reaches only $89/90$, because $\NK(r)=r^2$ leaves only $Y^{1/2}$ rows.
\item \textbf{The prime count needed is now available.} The 7/8 chain's prime ideal count
(\lean{AnnularPrimeMass}) would discharge \lean{FixedTestPrimeMS}, but that would add fidelity to
Paper~2, not a better result.
\end{itemize}

\subsection{Descent, inversion and dictionary (7/8)}
There is no hidden alternative proof here: all 914 dead declarations are finished. The largest group,
about $2.7$k lines, is a ladder of re-proofs forced by the shape of the final induction predicate:
\begin{itemize}
\item windows must be chosen before the degree;
\item constants must be uniform in the index types;
\item the degree must be uniform in the level.
\end{itemize}
This is the 7/8 paper's lemma on finite seminorm propagation. It is stated abstractly there, and the
formalization rediscovered it only when the top-level induction failed.

Other replacements:
\begin{itemize}
\item Path-length and real-parameter inductions became induction on an integer rank. The formal
losses per step are $15\eta$ and $62\eta$, against the paper's $11\eta$ and $40\eta+\tau+\pi$; this
does not affect the exponent.
\item Per-source radii and slot lists became data that is constant on each cell.
\item Sums over one fixed cube became a global source, keeping the cube sum inside the square.
\item Pools that depend on the tuples became the excluded-prime pool.
\end{itemize}

No variant implements Paper~2's adaptive cube cutoff $H_c$, and none gives a better exponent.

\subsection{Fourth moments (7/8, \texorpdfstring{\lean{Moments/}}{Moments/})}
\begin{itemize}
\item \textbf{The rectangle cancellation needed explicit seminorm constants.} Its constants went
through four forms of uniformity before reaching explicit seminorm products. The form ``$\exists J,
\forall W, \exists C$'' cannot feed an induction whose test functions move.
\item \textbf{A decentered majorant was a dead end.} One chain majorized the rectangle difference by
$|S(X)|^2+|S(Y)|^2$. That gives up the $\Theta$ cancellation the centered range needs, although it
is valid for the uncentered range $A\le5M/6$.
\item \textbf{The descent hypotheses were tuned rather than proved.} The second-transform descent's
hypotheses on child energies could not be met at the recursive call. The live versions do two
things:
\begin{itemize}
\item restrict to children that actually occur, the ``frozen'' variant, at the cost of a factor
$(n_2n_3/n_0n_1)^{1/6}$;
\item pay $1/\NK(\mathrm{seed})$ explicitly in the tail term, the ``seeded'' variant.
\end{itemize}
\item \textbf{One paragraph of the 7/8 paper was formalized but is unnecessary.} The 7/8 paper's
``added row zero'' estimate (sixth powers among the rows) is formalized verbatim. The formal proof
keeps the zero row punctured, so it never needs it.
\item \textbf{There is no slack in the exponent bookkeeping.} Dead ledger lemmas certify the paper's
exponent inequality, and show it is tight at $v=M/4$.
\end{itemize}

\subsection{Detector, energy and the 7/8 assembly}
Most of this is the ``before'' version of three changes in the order of choices:
\begin{itemize}
\item a hypothesis ``for every slot datum $D$'' became ``for some fine $D$'', as in the 7/8 paper's
proposition on the order of choices;
\item constants chosen inside a theorem became an input structure;
\item a fixed stage loss became a parameter, together with a cap on the level.
\end{itemize}
The assembly over every zero of every character was replaced by one over primitive characters only.
The uniform contour estimates were stronger than the paper requires.

\subsection{Sieves, Gauss sums and the mean square}
\begin{itemize}
\item \textbf{One conditional chain.} Version~1 of the initial dual-mean-square density,
\lean{HasInitialCanonicalDensity}, is assumed but never constructed. It demanded one derivative
order for all test windows. The descent doubles the window's support radius at every level ($A\mapsto
2A+11$), so the live version \lean{...AtRadius} lets that order depend on the radius.
\item \textbf{Two-passage transfers.} The largest cluster, about $7$k lines, is a family of
monolithic ``first-plus-second Poisson'' transfer inequalities with bespoke budget structures. They
were superseded by a modular split: the first passage produces an abstract output energy, the second
passage gives twisted children, and a closing step finishes. \emph{Conjecture:} the monolithic
outputs could not be fed back into the recursion without re-binning.
\item \textbf{Quadratic large sieve.} No Goldmakher--Louvel attempt exists. Earlier Heath-Brown-style
recursions on block norms defined by suprema stopped at a recursive inequality whose block norms
were never bounded. The live proof bootstraps a single scalar exponent, $\alpha\mapsto2-1/\alpha$.
\item \textbf{Uniformity again.} The non-uniform versions of Prop.~R and of the transfers each fail
exactly one uniformity that the descent requires: uniformity in the prime pool, in the height twist,
or in the window family.
\item \textbf{Absent from the 7/8 directories too.} There is no band-limited weight $\Phi$, no
analogue of Paper~2's lemma removing excluded primes, and no descent with a general gap or cube
cutoff. The numbers $1/20$, $1/1000$ and $3000$ are fixed everywhere.
\item \textbf{Useful leftovers.}
\begin{itemize}
\item Two complete, constant-free descent lemmas: \lean{ShortDraftDescent.integral_geometric_descent}
and \lean{binary_energy_descent}.
\item A $\sigma$-general initial invariant, \lean{initial_invariant_power_general}.
\item The paper-faithful prime-row extraction, whose exponents are linear in $\vartheta$.
\end{itemize}
\end{itemize}

\section{Ideas worth more compute}\label{sec:leads}
Ranked by expected value per unit of effort.
\begin{enumerate}
\item \textbf{A 15/16 version of the short proof (cheap; verified prerequisite).}
\begin{itemize}
\item The formal mean-square bound
(\lean{InitialMeanSquare.outside_mean_square_of_canonical_density_at_radius}) holds for every row
range $D^{1+\sigma}$ with $\sigma>1/20$, and its other input does not depend on $\sigma$. Only the
extraction and the Mellin step hard-code $\vartheta=1/10$, through the literals
$11/60, 11/10, 21/10, 137/120, 23/24$.
\item Making them parametric gives a zero-free half-plane $\Rea s>11/12+5\vartheta/12$ for each
$\vartheta>1/20$, hence $\Rea s>15/16$ after letting $\vartheta$ tend to $1/20$.
\item The estimated cost is 100--150 Lean lines in the top-level files.
\item Reaching $11/12$ itself would further require making the descent's fixed margin $1/20$ a
parameter (Remark~5.5 of the 23/24 paper).
\end{itemize}
\item \textbf{A parameterised descent, towards $11/12$.}
\begin{itemize}
\item Rebuild the 23/24 descent on the dead, constant-free
\lean{ShortDraftDescent.integral_geometric_descent}. Its rank should be
$\lceil\log K/(c\log Z)\rceil$, with mass $C_rZ^{\varepsilon_r}(XF)^2$.
\item Seed it with \lean{initial_invariant_power_general}, which already holds for every $\sigma$.
\item Thread the gap through the few lemmas where $1/20$, $1/40$, $1/1000$ and $3000$ enter.
\end{itemize}
Together with item~1, this is the formal route to $11/12$. Its cost has not been measured.
\item \textbf{Remove the exclusion mask from the descent state (a cheap test).}
\begin{itemize}
\item The live descent state carries a mask $m$ that grows down the tree, and its invariant charges
$\|m\|^2$ against the gap.
\item A Lean analogue of Paper~2's lemma removing excluded primes,
\[E_m\le\tau(m)\sum_{d\mid m}E(K,X/\NK d,F\NK d),\]
starting from the inclusion--exclusion identity
\lean{ideal_mobius_coprime_reindex}, would remove the mask.
\item That might shorten the descent. \emph{Unverified.}
\end{itemize}
\item \textbf{Close Paper~2's $\mathrm{SL}_2(\mathbb Z)$-invariance gap.}
\begin{itemize}
\item Live and dead lemmas together reduce ``$\theta$ equals its average'' to one finite identity
between two cubic Gauss-sum residues for horizontal $h$:
\[
\tfrac{9\sqrt3}{2}\cdot\text{\lean{principalArithmeticResidue}}(h)=\text{\lean{unramifiedGaussResidue}}(9h).
\]
The constant terms already agree, and both sides vanish at non-horizontal $h$.
\item Proving the identity would make the formal $\theta$ literally Kubota's, matching Paper~2
without averaging.
\item The ``decaying eigenfunction $\Rightarrow0$'' template already used for $\tau(hp^3)$ is an
alternative route.
\item If the identity \emph{fails}, that is itself a finding about Paper~2.
\end{itemize}
\item \textbf{Shrink the theta construction}, about 30\% of the 23/24 proof's code.
\begin{itemize}
\item The dead code shows that meromorphy, self-adjointness, non-vanishing of the residue and
finite-dimensionality are all unused.
\item Package the Fredholm pencil and its Laurent split as one generic ``simple-pole residue''
lemma, then re-run the reachability audit.
\end{itemize}
\item \textbf{Shorten the 7/8 proof.}
\begin{itemize}
\item Bundle the roughly 40 repeated hypotheses of the surviving descent ladder: theorem statements
are 52\% of its 2,191 lines.
\item Try the short, generic hybrid quadratic--cubic sieve lemmas in place of about 1,150 lines of
grid machinery.
\item Use the dead decentered majorant together with \lean{actual_live_slot_removal} for the
uncentered range $A\le5M/6$.
\end{itemize}
\item \textbf{Smaller items.}
\begin{itemize}
\item Run a linear-programming slack analysis of the 7/8 exponent constraints, to find which ones
bind. The known one, at $M/4$, is tight.
\item Drop the added-zero paragraph from the 7/8 paper.
\item Upstream the self-contained Mellin criteria to Mathlib.
\end{itemize}
\end{enumerate}

\section{Dead ends worth avoiding}\label{sec:deadends}
\subsection*{Mathematics}
\begin{itemize}
\item Reflection with $\Gamma(3)$-automorphy alone: the Weyl-type cusps need $\mathrm{SL}_2(\mathbb
Z)$. Average first, or prove the invariance first.
\item Global spectral theory, or Green and Maass--Selberg identities on cusp boxes in hyperbolic
3-space: a local Laurent split at $s=4/3$ suffices.
\item Euler-factor recurrences for $\tau(hp^3)$: they leave one constant free.
\item Cauchy--Schwarz over all $3^{|\varphi|}$ cusp branches: restrict to the surviving branches
first.
\item Large-sieve recursions on block norms defined by suprema: the scalar exponent bootstrap
closes, the block-norm recursion never did.
\item One monolithic two-Poisson transfer inequality with a bespoke budget: keep the passes modular.
\item Decentering a rectangle difference before the $\Theta$ cancellation.
\item Prime rows instead of integer rows: they give the same exponent, and need a prime ideal
theorem. Rational-prime rows are strictly worse ($89/90$).
\item Re-tuning the 7/8 exponent ledgers: they close with zero slack.
\end{itemize}
\subsection*{Process (for agent-driven formalization)}
\begin{itemize}
\item \textbf{Fix the consumer before the producer.} Interfaces written before their inputs existed
were rewritten six times (the Mellin step). Stage lemmas proved before the induction that consumes
them were rewritten as well.
\item \textbf{Settle quantifier order and uniformity first.} Most dead lines are the same theorem
with $\exists$ and $\forall$ in a different order. Fix the shape of the top-level induction
predicate first.
\item \textbf{Check that a hypothesis can be discharged at its call site before building a chain on
it.} A hypothesis-carrying structure was a warning sign every time.
\item \textbf{Generalize by adding a parameter, not by copying.} Clones were usually imported by
their successors, so they were also compiled on every build.
\item \textbf{Prove only the uniformity the paper needs, and match its representation.} Punctured
rows made one formalized paragraph of the paper unnecessary.
\end{itemize}

\end{document}
